% ch05_structures.tex -- Data Structures
\chapter{Data Structures}
\label{ch:structures}

\begin{intuition}
A scientist never works with a single data point: they handle series of
measurements, tables of results, catalogs of species. Python offers several
\emph{data structures} to organize this information --- each with its strengths,
like the different containers in a laboratory (test tubes, Petri dishes,
binders).
\end{intuition}

% ============================================================
\section{Lists}
\index{list}

\begin{definition}[List]
A \textbf{list} is an \emph{ordered} and \emph{mutable} collection of elements.
It is created with square brackets: \py{[e1, e2, ...]}.
\end{definition}

\subsection{Creation and indexing}
\index{indexing}

\begin{codeblock}[title=\textbf{Python}]
\begin{minted}{python}
# Temperature measurements (deg C) over one week
temperatures = [18.2, 19.5, 17.8, 21.0, 22.3, 20.1, 19.7]

print(f"First day    : {temperatures[0]} C")
print(f"Last day     : {temperatures[-1]} C")
print(f"Number of measurements: {len(temperatures)}")
\end{minted}
\end{codeblock}

\begin{codeoutput}
\begin{verbatim}
First day    : 18.2 C
Last day     : 19.7 C
Number of measurements: 7
\end{verbatim}
\end{codeoutput}

\subsection{Slicing}
\index{slicing}

\begin{codeblock}[title=\textbf{Python}]
\begin{minted}{python}
# Weekdays (Monday to Friday)
jours_ouvres = temperatures[0:5]
print(f"Weekdays: {jours_ouvres}")

# Every other element
print(f"Every other: {temperatures[::2]}")

# Reverse order
print(f"Reversed: {temperatures[::-1]}")
\end{minted}
\end{codeblock}

\begin{codeoutput}
\begin{verbatim}
Weekdays: [18.2, 19.5, 17.8, 21.0, 22.3]
Every other: [18.2, 17.8, 22.3, 19.7]
Reversed: [19.7, 20.1, 22.3, 21.0, 17.8, 19.5, 18.2]
\end{verbatim}
\end{codeoutput}

\subsection{Useful methods}
\index{append}\index{pop}\index{sort}

\begin{codeblock}[title=\textbf{Python}]
\begin{minted}{python}
masses = [12.5, 8.3, 15.7, 10.2]

masses.append(9.8)       # add to the end
print(f"After append: {masses}")

dernier = masses.pop()   # remove the last element
print(f"Removed: {dernier}, list: {masses}")

masses.sort()            # sort in place
print(f"Sorted: {masses}")

masses.reverse()         # reverse in place
print(f"Reversed: {masses}")
\end{minted}
\end{codeblock}

\begin{codeoutput}
\begin{verbatim}
After append: [12.5, 8.3, 15.7, 10.2, 9.8]
Removed: 9.8, list: [12.5, 8.3, 15.7, 10.2]
Sorted: [8.3, 10.2, 12.5, 15.7]
Reversed: [15.7, 12.5, 10.2, 8.3]
\end{verbatim}
\end{codeoutput}

% ============================================================
\section{List comprehensions}
\index{list comprehension}

List comprehensions offer a compact syntax for creating lists.

\begin{codeblock}[title=\textbf{Python}]
\begin{minted}{python}
# Squares of integers from 1 to 10
carres = [x**2 for x in range(1, 11)]
print(f"Squares: {carres}")

# Filter temperatures > 20 deg C
temperatures = [18.2, 19.5, 17.8, 21.0, 22.3, 20.1, 19.7]
chaudes = [t for t in temperatures if t > 20]
print(f"Hot days: {chaudes}")

# Convert to Fahrenheit
fahrenheit = [t * 9/5 + 32 for t in temperatures]
print(f"Fahrenheit: {[f'{f:.1f}' for f in fahrenheit]}")
\end{minted}
\end{codeblock}

\begin{codeoutput}
\begin{verbatim}
Squares: [1, 4, 9, 16, 25, 36, 49, 64, 81, 100]
Hot days: [21.0, 22.3, 20.1]
Fahrenheit: ['64.8', '67.1', '64.0', '69.8', '72.1', '68.2', '67.5']
\end{verbatim}
\end{codeoutput}

% ============================================================
\section{Tuples}
\index{tuple}

\begin{definition}[Tuple]
A \textbf{tuple} is an \emph{ordered} but \emph{immutable}
(non-modifiable) collection. It is created with parentheses: \py{(e1, e2, ...)}.
\end{definition}

\begin{codeblock}[title=\textbf{Python}]
\begin{minted}{python}
# Physical constants (must not change)
constantes = ("speed of light", 299_792_458, "m/s")
print(f"{constantes[0]} = {constantes[1]} {constantes[2]}")

# Unpacking
nom, valeur, unite = constantes
print(f"{nom}: {valeur} {unite}")

# Attempting to modify -> error
# constantes[1] = 300_000_000  # TypeError!
\end{minted}
\end{codeblock}

\begin{codeoutput}
\begin{verbatim}
speed of light = 299792458 m/s
speed of light: 299792458 m/s
\end{verbatim}
\end{codeoutput}

\begin{remark}
Use tuples for data that should not be modified:
coordinates $(x, y)$, physical constants, dictionary keys.
\end{remark}

% ============================================================
\section{Dictionaries}
\index{dictionary}

\begin{definition}[Dictionary]
A \textbf{dictionary} maps \emph{keys} to \emph{values}.
It is created with curly braces: \py{\{key: value, ...\}}.
\end{definition}

\subsection{The periodic table as a dictionary}
\index{periodic table}

\begin{codeblock}[title=\textbf{Python}]
\begin{minted}{python}
elements = {
    "H":  {"name": "Hydrogen",  "Z": 1,  "mass": 1.008},
    "He": {"name": "Helium",    "Z": 2,  "mass": 4.003},
    "C":  {"name": "Carbon",    "Z": 6,  "mass": 12.011},
    "N":  {"name": "Nitrogen",  "Z": 7,  "mass": 14.007},
    "O":  {"name": "Oxygen",    "Z": 8,  "mass": 15.999},
    "Fe": {"name": "Iron",      "Z": 26, "mass": 55.845},
}

# Access an element
carbone = elements["C"]
print(f"Carbon: Z={carbone['Z']}, mass={carbone['mass']} u")

# Molar mass of water H2O
masse_eau = 2 * elements["H"]["mass"] + elements["O"]["mass"]
print(f"Molar mass H2O = {masse_eau:.3f} g/mol")
\end{minted}
\end{codeblock}

\begin{codeoutput}
\begin{verbatim}
Carbon: Z=6, mass=12.011 u
Molar mass H2O = 18.015 g/mol
\end{verbatim}
\end{codeoutput}

\subsection{Iterating over a dictionary}

\begin{codeblock}[title=\textbf{Python}]
\begin{minted}{python}
# Display all elements
for symbole, info in elements.items():
    print(f"{symbole:>2} : {info['name']:<12} Z={info['Z']:>2}")
\end{minted}
\end{codeblock}

\begin{codeoutput}
\begin{verbatim}
 H : Hydrogen     Z= 1
He : Helium       Z= 2
 C : Carbon       Z= 6
 N : Nitrogen     Z= 7
 O : Oxygen       Z= 8
Fe : Iron         Z=26
\end{verbatim}
\end{codeoutput}

% ============================================================
\section{Sets}
\index{set}

\begin{codeblock}[title=\textbf{Python}]
\begin{minted}{python}
# Elements detected in two samples
echantillon_A = {"H", "C", "O", "N", "S"}
echantillon_B = {"H", "C", "O", "Fe", "Ca"}

print(f"Common      : {echantillon_A & echantillon_B}")
print(f"All         : {echantillon_A | echantillon_B}")
print(f"Only in A   : {echantillon_A - echantillon_B}")
\end{minted}
\end{codeblock}

\begin{codeoutput}
\begin{verbatim}
Common      : {'O', 'H', 'C'}
All         : {'S', 'O', 'Fe', 'H', 'Ca', 'N', 'C'}
Only in A   : {'S', 'N'}
\end{verbatim}
\end{codeoutput}

% ============================================================
% TikZ comparison
\begin{figure}[ht]
\centering
\begin{tikzpicture}[
    header/.style={rectangle, draw, fill=blue!60, text=white,
        minimum width=3.5cm, minimum height=0.7cm, font=\bfseries, align=center},
    cell/.style={rectangle, draw, minimum width=3.5cm, minimum height=0.6cm,
        font=\small, align=center},
]
% List
\node[header] (lh) {List};
\node[cell, below=0pt of lh] (l1) {Ordered};
\node[cell, below=0pt of l1] (l2) {Mutable};
\node[cell, below=0pt of l2] (l3) {\py{[1, 2, 3]}};

% Tuple
\node[header, right=0.5cm of lh] (th) {Tuple};
\node[cell, below=0pt of th] (t1) {Ordered};
\node[cell, below=0pt of t1] (t2) {Immutable};
\node[cell, below=0pt of t2] (t3) {\py{(1, 2, 3)}};

% Dict
\node[header, right=0.5cm of th] (dh) {Dictionary};
\node[cell, below=0pt of dh] (d1) {Key $\to$ Value};
\node[cell, below=0pt of d1] (d2) {Mutable};
\node[cell, below=0pt of d2] (d3) {\py{\{"a": 1\}}};

% Set
\node[header, right=0.5cm of dh] (sh) {Set};
\node[cell, below=0pt of sh] (s1) {Unordered};
\node[cell, below=0pt of s1] (s2) {No duplicates};
\node[cell, below=0pt of s2] (s3) {\py{\{1, 2, 3\}}};
\end{tikzpicture}
\caption{Comparison of data structures in Python.}
\label{fig:structures-comparaison}
\end{figure}

% ============================================================
\section{Common errors}

\begin{errmark}
\begin{enumerate}
    \item \textbf{Index out of bounds (\py{IndexError})}: a list with 5 elements has
    indices from 0 to 4. Accessing index 5 causes an error.
    \begin{minted}{python}
mesures = [1.2, 3.4, 5.6]
print(mesures[3])  # IndexError!
# Correct: mesures[2] or mesures[-1]
    \end{minted}

    \item \textbf{Modifying a list during iteration}: removing elements
    from a list while iterating over it produces unpredictable results.
    \begin{minted}{python}
# ERROR: skips elements
donnees = [1, -2, 3, -4, 5]
for d in donnees:
    if d < 0:
        donnees.remove(d)  # Dangerous!

# CORRECT: create a new list
donnees = [d for d in donnees if d >= 0]
    \end{minted}

    \item \textbf{Missing key (\py{KeyError})}: accessing a key that does not exist in a
    dictionary. Use \py{dict.get(key, default)} to avoid the error.

    \item \textbf{Confusing list and tuple}: attempting to modify a tuple with
    \py{t[0] = 5} causes a \py{TypeError}.
\end{enumerate}
\end{errmark}

% ============================================================
\section{Mini-project: Student Grade Management}

\begin{projet}
Create a grade management system using a dictionary:
\begin{itemize}
    \item Keys: student names
    \item Values: lists of grades
    \item Functions: add a student, add a grade, compute the average,
          find the top student
\end{itemize}
\end{projet}

\begin{codeblock}[title=\textbf{Python}]
\begin{minted}{python}
def creer_promotion():
    """Create an empty class dictionary."""
    return {}

def ajouter_etudiant(promo, nom):
    """Add a student with an empty grade list."""
    promo[nom] = []

def ajouter_note(promo, nom, note):
    """Add a grade for a student."""
    promo[nom].append(note)

def moyenne(notes):
    """Compute the average of a list of grades."""
    return sum(notes) / len(notes) if notes else 0

def afficher_resultats(promo):
    """Display the results of the class."""
    print(f"{'Student':<15} {'Grades':<25} {'Average':>8}")
    print("-" * 50)
    for nom, notes in promo.items():
        notes_str = ", ".join(f"{n:.1f}" for n in notes)
        moy = moyenne(notes)
        print(f"{nom:<15} {notes_str:<25} {moy:>7.2f}")

# Usage
promo = creer_promotion()
for nom in ["Alice", "Bob", "Claire"]:
    ajouter_etudiant(promo, nom)

ajouter_note(promo, "Alice", 15.5)
ajouter_note(promo, "Alice", 17.0)
ajouter_note(promo, "Alice", 14.0)
ajouter_note(promo, "Bob", 12.0)
ajouter_note(promo, "Bob", 11.5)
ajouter_note(promo, "Bob", 13.0)
ajouter_note(promo, "Claire", 18.0)
ajouter_note(promo, "Claire", 16.5)
ajouter_note(promo, "Claire", 19.0)

afficher_resultats(promo)

major = max(promo, key=lambda nom: moyenne(promo[nom]))
print(f"\nTop student: {major} ({moyenne(promo[major]):.2f}/20)")
\end{minted}
\end{codeblock}

\begin{codeoutput}
\begin{verbatim}
Student         Grades                    Average
--------------------------------------------------
Alice           15.5, 17.0, 14.0            15.50
Bob             12.0, 11.5, 13.0            12.17
Claire          18.0, 16.5, 19.0            17.83

Top student: Claire (17.83/20)
\end{verbatim}
\end{codeoutput}

% ============================================================
\section{Exercises}

\begin{exercise}[$\star$ --- List Manipulation]
Create a list of 10 pH measurements. Compute the mean, minimum, and maximum
\emph{without using} the \py{sum}, \py{min}, \py{max} functions.
\end{exercise}

\begin{exercise}[$\star$ --- Filtering]
From a list of temperatures, create a new list using a list comprehension
containing only values between 15\,\textdegree C and 25\,\textdegree C.
\end{exercise}

\begin{exercise}[$\star\star$ --- Frequency Counter]
Write a function that receives a string and returns a dictionary
counting the frequency of each letter. Test with a DNA sequence
(``ATCGGATCCA'').
\end{exercise}

\begin{exercise}[$\star\star$ --- Matrix as a List of Lists]
Represent a $3\times 3$ matrix as a list of lists. Write a function
that computes the transpose. Verify with the identity matrix.
\end{exercise}

\begin{exercise}[$\star\star\star$ --- Protein Sequence Analysis]
Create a dictionary mapping each amino acid to its molar mass.
Write a function that, given a sequence of amino acids (letters),
computes the total mass of the protein. Test with ``MKVLWA''.
\end{exercise}

% ============================================================
\begin{keyformulas}{Chapter Summary: Data Structures}
\begin{itemize}
    \item \textbf{List} \py{[...]}: ordered, mutable, methods \py{append},
    \py{pop}, \py{sort}
    \item \textbf{Comprehension}: \py{[expr for x in iterable if cond]}
    \item \textbf{Tuple} \py{(...)}: ordered, immutable, ideal for constants
    \item \textbf{Dictionary} \py{\{key: val\}}: access by key, methods
    \py{.keys()}, \py{.values()}, \py{.items()}
    \item \textbf{Set} \py{\{...\}}: no duplicates, operations $\cap$, $\cup$, $\setminus$
    \item \textbf{Indexing}: starts at 0, negative indices from the end
    \item \textbf{Slicing}: \py{list[start:stop:step]}
\end{itemize}
\end{keyformulas}
